\chapter{Prefácio}

O autor é o primeiro a admitir que o presente material representa uma aproximação inicial de unidade curricular inédita. Assim, é possível (e esperado) que ele sofra substanciais alterações entre as ofertas. Isso não é uma defesa acerca da qualidade do conteúdo --- produzido com muito esmero --- mas uma aquiescência de que a execução da disciplina trará experiências que só poderão enriquecer ofertas futuras.

A unidade curricular IAS1003 (Fundamentos de Modelagem e Análise de Dados) não ocorre num vácuo. Ela é apenas uma de outras 6 unidades curriculares do primeiro período do Curso Superior de Tecnologia em Inteligência Artificial e Sistemas Autônomos. Há uma distinta vantagem de lecionar no primeiro período: todos os discentes estão matriculados no mesmo conjunto de unidades curriculares, sendo elas:
\begin{itemize}
	\item Fundamentos de Eletricidade e Eletrônica;
	\item Sistemas Digitais;
	\item \emph{Fundamentos de Modelagem e Análise de Dados};
	\item Programação para Sistemas Autônomos;
	\item Fábricas Inteligentes e Fundamentos de Autonomia;
	\item Ética, Sustentabilidade e Humanização da Tecnologia.
\end{itemize}
